\section{Cross-Border Flows and Market Coupling}

\subsection{Why PL price is not an island}

The Polish bidding zone is connected to six neighbours:
Germany/Luxemburg (DE-LU), Czechia (CZ), Slovakia (SK),
Lithuania (LT), Sweden (SE4), and Ukraine (UA-BEI, limited).

When German price is lower than Polish price, power flows in (import).
The physical flow equals the scheduled cross-border exchange.
The flow continues until the prices converge (arbitrage closes the gap)
or the interconnector hits its capacity limit.

That limit is the Net Transfer Capacity (NTC) or, after implementation,
the Flow-Based Capacity (FBC) model.

\subsection{Market coupling: SDAC}

Poland participates in the Single Day-Ahead Coupling (SDAC),
which optimises allocation of cross-border capacity and market matching
simultaneously across the whole CWE/CEE region.

The SDAC algorithm (EUPHEMIA) does not allocate capacity separately
from the energy auction — it co-optimises. This means:

\begin{itemize}
  \item Interconnector capacity is allocated to whoever values it most.
  \item Prices tend to converge across coupled zones.
  \item Isolated zones (Ukraine, until recently Belarus) have no coupling
        benefit — they face their own merit order.
\end{itemize}

The result: PL day-ahead price is increasingly correlated with DE price
as inter-connector capacity grows. A model that ignores cross-border
flows misses this coupling signal.

\subsection{Price convergence and congestion}

When lines are not congested: PL price = DE price + transmission losses (small).
When lines are congested: prices diverge. PL can be higher than DE
even with a surplus in Germany if the cable is full.

In practice, PL is a net importer in winter (more demand than domestic capacity)
and a net exporter when domestic solar + wind peaks.
Congestion is common in both directions during extreme events.

Measured in our data: PL--DE price spread averages 4.2 EUR/MWh (absolute).
On congested hours, spread can exceed 50 EUR/MWh.

\subsection{Net position as a feature}

The scheduled net position (net export/import for PL) is published by ENTSO-E.
It is available for day D at roughly 13:00 on D-1 (after SDAC results).
This is after our 09:00 cutoff.

A proxy available at cutoff: the TSO RES forecast implies the residual
load, which implies whether PL will be exporting or importing.
This is why the TSO load and RES forecasts are already strong proxies
for the cross-border signal.

Adding neighbour prices as a direct lag feature (lag-24h DE price,
lag-24h CZ price) is the next upgrade. It would capture:
- Periods when PL is tightly coupled to a cheap-price neighbour
- Gas crisis propagation from the TTF-exposed DE market

This feature addition requires config/one-file change only
(modular input design, see CLAUDE.md). Implementation time: \textless 1h.

\subsection{Flow-Based Capacity (FBC): the structural change}

ENTSO-E is migrating the CWE region from NTC (net transfer capacity —
a simple per-interconnector limit) to Flow-Based Capacity modelling.
Under FBC, the available capacity depends on the actual power flows
through the transmission grid (the Power Transfer Distribution Factor,
PTDF, matrix).

FBC affects how quickly prices converge across borders. A model trained
on NTC data may mis-estimate correlation under FBC.
Poland is scheduled to join the FBC region, but the timeline has slipped.

\subsection{Interview lines}

\begin{itemize}
  \item ``PL price is market-coupled to six neighbours via SDAC. When
        interconnectors are not congested, prices converge. Our model
        captures this through the TSO RES forecast proxy, which predicts
        net position direction. A direct lag of DE-LU price is the next
        feature on the backlog.''
  \item ``PL--DE price spread averages 4.2 EUR/MWh but exceeds 50 EUR/MWh
        on congested hours. Congestion events align with extreme weather
        and should be captured by the merit-order feature set.''
  \item ``I did not add neighbour prices in Phase 2 because the modular
        feature design (config-driven, one file per concern) means it is
        a single pipeline change. I deferred it to Phase 3 to avoid
        scope creep during the shadow window.''
\end{itemize}
